\documentclass[14pt, russian]{scrartcl}
\let\counterwithout\relax
\let\counterwithin\relax
\usepackage{float}
\usepackage{xcolor}
\usepackage{extsizes}
\usepackage{subfig}
\usepackage[export]{adjustbox}
\usepackage{tocvsec2} % возможность менять учитываемую глубину разделов в оглавлении
\usepackage[subfigure]{tocloft}
\usepackage[newfloat]{minted}
\captionsetup[listing]{position=top}

\makeatletter
\AddToHook{begindocument/before}{\@ifpackageloaded{minted}{\removefromtoclist[float]{lol}}{}}
\makeatother

\AtBeginEnvironment{figure}{\vspace{0.5cm}}
\AtBeginEnvironment{table}{\vspace{0.5cm}}
\AtBeginEnvironment{listing}{\vspace{0.5cm}}
\AtBeginEnvironment{algorithm}{\vspace{0.5cm}}
\AtBeginEnvironment{minted}{\vspace{-0.5cm}}

\usepackage{fancyvrb}
\usepackage{ulem,bm,mathrsfs,ifsym} %зачеркивания, особо жирный стиль и RSFS начертание
\usepackage{sectsty} % переопределение стилей подразделов
%%%%%%%%%%%%%%%%%%%%%%%

%%% Поля и разметка страницы %%%
\usepackage{pdflscape}                              % Для включения альбомных страниц
\usepackage{geometry}                               % Для последующего задания полей
\geometry{a4paper,tmargin=2cm,bmargin=2cm,lmargin=3cm,rmargin=1cm} % тоже самое, но лучше

%%% Математические пакеты %%%
\usepackage{amsthm,amsfonts,amsmath,amssymb,amscd}  % Математические дополнения от AMS
\usepackage{mathtools}                              % Добавляет окружение multlined
\usepackage[perpage]{footmisc}
%\usepackage{times}

%%%% Установки для размера шрифта 14 pt %%%%
%% Формирование переменных и констант для сравнения (один раз для всех подключаемых файлов)%%
%% должно располагаться до вызова пакета fontspec или polyglossia, потому что они сбивают его работу
%\newlength{\curtextsize}
%\newlength{\bigtextsize}
%\setlength{\bigtextsize}{13pt}
\KOMAoptions{fontsize=14pt}

\makeatletter
\def\showfontsize{\f@size{} point}
\makeatother

\makeatletter
\setlength{\@fptop}{0pt}
\makeatother

%\makeatletter
%\show\f@size                                       % неплохо для отслеживания, но вызывает стопорение процесса, если документ компилируется без команды  -interaction=nonstopmode 
%\setlength{\curtextsize}{\f@size pt}
%\makeatother

%шрифт times
\usepackage{tempora}
%\usepackage{pscyr}
%\setmainfont[Ligatures={TeX,Historic}]{Times New Roman}

   %%% Решение проблемы копирования текста в буфер кракозябрами
%    \input glyphtounicode.tex
%    \input glyphtounicode-cmr.tex %from pdfx package
%    \pdfgentounicode=1
    \usepackage{cmap}                               % Улучшенный поиск русских слов в полученном pdf-файле
    \usepackage[T1]{fontenc}                       % Поддержка русских букв
    \usepackage[utf8]{inputenc}                     % Кодировка utf8
    \usepackage[english, main=russian]{babel}            % Языки: русский, английский
%   \IfFileExists{pscyr.sty}{\usepackage{pscyr}}{}  % Красивые русские шрифты
%\renewcommand{\rmdefault}{ftm}
%%% Оформление абзацев %%%
\usepackage{indentfirst}                            % Красная строка
%\usepackage{eskdpz}

%%% Таблицы %%%
\usepackage{longtable}                              % Длинные таблицы
\usepackage{multirow,makecell,array}                % Улучшенное форматирование таблиц
\usepackage{booktabs}                               % Возможность оформления таблиц в классическом книжном стиле (при правильном использовании не противоречит ГОСТ)

%%% Общее форматирование
\usepackage{soulutf8}                               % Поддержка переносоустойчивых подчёркиваний и зачёркиваний
\usepackage{icomma}                                 % Запятая в десятичных дробях



%%% Изображения %%%
\usepackage{graphicx}                               % Подключаем пакет работы с графикой
\usepackage{wrapfig}

%%% Списки %%%
\usepackage{enumitem}

%%% Подписи %%%
\usepackage{caption}                                % Для управления подписями (рисунков и таблиц) % Может управлять номерами рисунков и таблиц с caption %Иногда может управлять заголовками в списках рисунков и таблиц
%% Использование:
%\begin{table}[h!]\ContinuedFloat - чтобы не переключать счетчик
%\captionsetup{labelformat=continued}% должен стоять до самого caption
%\caption{}
% либо ручками \caption*{Продолжение таблицы~\ref{...}.} :)

%%% Интервалы %%%
\addto\captionsrussian{%
  \renewcommand{\listingname}{Листинг}%
}
%%% Счётчики %%%
\usepackage[figure,table,section]{totalcount}               % Счётчик рисунков и таблиц
\DeclareTotalCounter{lstlisting}
\usepackage{totcount}                               % Пакет создания счётчиков на основе последнего номера подсчитываемого элемента (может требовать дважды компилировать документ)
\usepackage{totpages}                               % Счётчик страниц, совместимый с hyperref (ссылается на номер последней страницы). Желательно ставить последним пакетом в преамбуле

%%% Продвинутое управление групповыми ссылками (пока только формулами) %%%
%% Кодировки и шрифты %%%

%   \newfontfamily{\cyrillicfont}{Times New Roman}
%   \newfontfamily{\cyrillicfonttt}{CMU Typewriter Text}
	%\setmainfont{Times New Roman}
	%\newfontfamily\cyrillicfont{Times New Roman}
	%\setsansfont{Times New Roman}                    %% задаёт шрифт без засечек
%	\setmonofont{Liberation Mono}               %% задаёт моноширинный шрифт
%    \IfFileExists{pscyr.sty}{\renewcommand{\rmdefault}{ftm}}{}
%%% Интервалы %%%
%linespread-реализация ближе к реализации полуторного интервала в ворде.
%setspace реализация заточена под шрифты 10, 11, 12pt, под остальные кегли хуже, но всё же ближе к типографской классике. 
\linespread{1.3}                    % Полуторный интервал (ГОСТ Р 7.0.11-2011, 5.3.6)
%\renewcommand{\@biblabel}[1]{#1}

%%% Гиперссылки %%%
\usepackage{hyperref}

%%% Выравнивание и переносы %%%
\sloppy                             % Избавляемся от переполнений
\clubpenalty=10000                  % Запрещаем разрыв страницы после первой строки абзаца
\widowpenalty=10000                 % Запрещаем разрыв страницы после последней строки абзаца

\makeatletter % малые заглавные, small caps shape
\let\@@scshape=\scshape
\renewcommand{\scshape}{%
  \ifnum\strcmp{\f@series}{bx}=\z@
    \usefont{T1}{cmr}{bx}{sc}%
  \else
    \ifnum\strcmp{\f@shape}{it}=\z@
      \fontshape{scsl}\selectfont
    \else
      \@@scshape
    \fi
  \fi}
\makeatother

%%% Подписи %%%
%\captionsetup{%
%singlelinecheck=off,                % Многострочные подписи, например у таблиц
%skip=2pt,                           % Вертикальная отбивка между подписью и содержимым рисунка или таблицы определяется ключом
%justification=centering,            % Центрирование подписей, заданных командой \caption
%}
%%%        Подключение пакетов                 %%%
\usepackage{ifthen}                 % добавляет ifthenelse
%%% Инициализирование переменных, не трогать!  %%%
\newcounter{intvl}
\newcounter{otstup}
\newcounter{contnumeq}
\newcounter{contnumfig}
\newcounter{contnumtab}
\newcounter{pgnum}
\newcounter{bibliosel}
\newcounter{chapstyle}
\newcounter{headingdelim}
\newcounter{headingalign}
\newcounter{headingsize}
\newcounter{tabcap}
\newcounter{tablaba}
\newcounter{tabtita}
%%%%%%%%%%%%%%%%%%%%%%%%%%%%%%%%%%%%%%%%%%%%%%%%%%

%%% Область упрощённого управления оформлением %%%

%% Интервал между заголовками и между заголовком и текстом
% Заголовки отделяют от текста сверху и снизу тремя интервалами (ГОСТ Р 7.0.11-2011, 5.3.5)
\setcounter{intvl}{3}               % Коэффициент кратности к размеру шрифта

%% Отступы у заголовков в тексте
\setcounter{otstup}{0}              % 0 --- без отступа; 1 --- абзацный отступ

%% Нумерация формул, таблиц и рисунков
\setcounter{contnumeq}{1}           % Нумерация формул: 0 --- пораздельно (во введении подряд, без номера раздела); 1 --- сквозная нумерация по всей диссертации
\setcounter{contnumfig}{1}          % Нумерация рисунков: 0 --- пораздельно (во введении подряд, без номера раздела); 1 --- сквозная нумерация по всей диссертации
\setcounter{contnumtab}{1}          % Нумерация таблиц: 0 --- пораздельно (во введении подряд, без номера раздела); 1 --- сквозная нумерация по всей диссертации

%% Оглавление
\setcounter{pgnum}{0}               % 0 --- номера страниц никак не обозначены; 1 --- Стр. над номерами страниц (дважды компилировать после изменения)

%% Библиография
\setcounter{bibliosel}{1}           % 0 --- встроенная реализация с загрузкой файла через движок bibtex8; 1 --- реализация пакетом biblatex через движок biber

%% Текст и форматирование заголовков
\setcounter{chapstyle}{1}           % 0 --- разделы только под номером; 1 --- разделы с названием "Глава" перед номером
\setcounter{headingdelim}{1}        % 0 --- номер отделен пропуском в 1em или \quad; 1 --- номера разделов и приложений отделены точкой с пробелом, подразделы пропуском без точки; 2 --- номера разделов, подразделов и приложений отделены точкой с пробелом.

%% Выравнивание заголовков в тексте
\setcounter{headingalign}{0}        % 0 --- по центру; 1 --- по левому краю

%% Размеры заголовков в тексте
\setcounter{headingsize}{0}         % 0 --- по ГОСТ, все всегда 14 пт; 1 --- пропорционально изменяющийся размер в зависимости от базового шрифта

%% Подпись таблиц
\setcounter{tabcap}{0}              % 0 --- по ГОСТ, номер таблицы и название разделены тире, выровнены по левому краю, при необходимости на нескольких строках; 1 --- подпись таблицы не по ГОСТ, на двух и более строках, дальнейшие настройки: 
%Выравнивание первой строки, с подписью и номером
\setcounter{tablaba}{2}             % 0 --- по левому краю; 1 --- по центру; 2 --- по правому краю
%Выравнивание строк с самим названием таблицы
\setcounter{tabtita}{1}             % 0 --- по левому краю; 1 --- по центру; 2 --- по правому краю

%%% Рисунки %%%
\DeclareCaptionLabelSeparator*{emdash}{~--- }             % (ГОСТ 2.105, 4.3.1)
\captionsetup[figure]{labelsep=emdash,font=onehalfspacing,position=bottom}

%%% Таблицы %%%
\ifthenelse{\equal{\thetabcap}{0}}{%
    \newcommand{\tabcapalign}{\raggedright}  % по левому краю страницы или аналога parbox
}

\ifthenelse{\equal{\thetablaba}{0} \AND \equal{\thetabcap}{1}}{%
    \newcommand{\tabcapalign}{\raggedright}  % по левому краю страницы или аналога parbox
}

\ifthenelse{\equal{\thetablaba}{1} \AND \equal{\thetabcap}{1}}{%
    \newcommand{\tabcapalign}{\centering}    % по центру страницы или аналога parbox
}

\ifthenelse{\equal{\thetablaba}{2} \AND \equal{\thetabcap}{1}}{%
    \newcommand{\tabcapalign}{\raggedleft}   % по правому краю страницы или аналога parbox
}

\ifthenelse{\equal{\thetabtita}{0} \AND \equal{\thetabcap}{1}}{%
    \newcommand{\tabtitalign}{\raggedright}  % по левому краю страницы или аналога parbox
}

\ifthenelse{\equal{\thetabtita}{1} \AND \equal{\thetabcap}{1}}{%
    \newcommand{\tabtitalign}{\centering}    % по центру страницы или аналога parbox
}

\ifthenelse{\equal{\thetabtita}{2} \AND \equal{\thetabcap}{1}}{%
    \newcommand{\tabtitalign}{\raggedleft}   % по правому краю страницы или аналога parbox
}

\DeclareCaptionFormat{tablenocaption}{\tabcapalign #1\strut}        % Наименование таблицы отсутствует
\ifthenelse{\equal{\thetabcap}{0}}{%
    \DeclareCaptionFormat{tablecaption}{\tabcapalign #1#2#3}
    \captionsetup[table]{labelsep=emdash}                       % тире как разделитель идентификатора с номером от наименования
}{%
    \DeclareCaptionFormat{tablecaption}{\tabcapalign #1#2\par%  % Идентификатор таблицы на отдельной строке
        \tabtitalign{#3}}                                       % Наименование таблицы строкой ниже
    \captionsetup[table]{labelsep=space}                        % пробельный разделитель идентификатора с номером от наименования
}
\captionsetup[table]{format=tablecaption,singlelinecheck=off,font=onehalfspacing,position=top,skip=-5pt}  % многострочные наименования и прочее
\DeclareCaptionLabelFormat{continued}{Продолжение таблицы~#2}
\setlength{\belowcaptionskip}{.2cm}
\setlength{\intextsep}{0ex}

%%% Подписи подрисунков %%%
\renewcommand{\thesubfigure}{\asbuk{subfigure}}           % Буквенные номера подрисунков
\captionsetup[subfigure]{font={normalsize},               % Шрифт подписи названий подрисунков (не отличается от основного)
    labelformat=brace,                                    % Формат обозначения подрисунка
    justification=centering,                              % Выключка подписей (форматирование), один из вариантов            
}
%\DeclareCaptionFont{font12pt}{\fontsize{12pt}{13pt}\selectfont} % объявляем шрифт 12pt для использования в подписях, тут же надо интерлиньяж объявлять, если не наследуется
%\captionsetup[subfigure]{font={font12pt}}                 % Шрифт подписи названий подрисунков (всегда 12pt)

%%% Настройки гиперссылок %%%

\definecolor{linkcolor}{rgb}{0.0,0,0}
\definecolor{citecolor}{rgb}{0,0.0,0}
\definecolor{urlcolor}{rgb}{0,0,0}

\hypersetup{
    linktocpage=true,           % ссылки с номера страницы в оглавлении, списке таблиц и списке рисунков
%    linktoc=all,                % both the section and page part are links
%    pdfpagelabels=false,        % set PDF page labels (true|false)
    plainpages=true,           % Forces page anchors to be named by the Arabic form  of the page number, rather than the formatted form
    colorlinks,                 % ссылки отображаются раскрашенным текстом, а не раскрашенным прямоугольником, вокруг текста
    linkcolor={linkcolor},      % цвет ссылок типа ref, eqref и подобных
    citecolor={citecolor},      % цвет ссылок-цитат
    urlcolor={urlcolor},        % цвет гиперссылок
    pdflang={ru},
}
\urlstyle{same}
%%% Шаблон %%%
%\DeclareRobustCommand{\todo}{\textcolor{red}}       % решаем проблему превращения названия цвета в результате \MakeUppercase, http://tex.stackexchange.com/a/187930/79756 , \DeclareRobustCommand protects \todo from expanding inside \MakeUppercase
\setlength{\parindent}{2.5em}                       % Абзацный отступ. Должен быть одинаковым по всему тексту и равен пяти знакам (ГОСТ Р 7.0.11-2011, 5.3.7).

%%% Списки %%%
% Используем дефис для ненумерованных списков (ГОСТ 2.105-95, 4.1.7)
%\renewcommand{\labelitemi}{\normalfont\bfseries~{---}} 
\renewcommand{\labelitemi}{\bfseries~{---}} 
\setlist{nosep,%                                    % Единый стиль для всех списков (пакет enumitem), без дополнительных интервалов.
    labelindent=\parindent,leftmargin=*%            % Каждый пункт, подпункт и перечисление записывают с абзацного отступа (ГОСТ 2.105-95, 4.1.8)
}
%%%%%%%%%%%%%%%%%%%%%%
%\usepackage{xltxtra} % load xunicode

\usepackage{ragged2e}
\usepackage[explicit]{titlesec}
\usepackage{placeins}
\usepackage{xparse}
\usepackage{csquotes}

\usepackage{listingsutf8}
\usepackage{url} %пакеты расширений
\usepackage{algorithm, algorithmicx}
\usepackage[noend]{algpseudocode}
\usepackage{blkarray}
\usepackage{chngcntr}
\usepackage{tabularx}
\usepackage[backend=biber, 
    bibstyle=gost-numeric,
    citestyle=nature]{biblatex}
\newcommand*\template[1]{\text{<}#1\text{>}}
  
\titleformat{name=\section,numberless}[block]{\normalfont\Large\centering}{}{0em}{#1}
\titleformat{\section}[block]{\normalfont\Large\bfseries\raggedright}{}{0em}{\thesection\hspace{0.25em}#1}
\titleformat{\subsection}[block]{\normalfont\Large\bfseries\raggedright}{}{0em}{\thesubsection\hspace{0.25em}#1}
\titleformat{\subsubsection}[block]{\normalfont\large\bfseries\raggedright}{}{0em}{\thesubsubsection\hspace{0.25em}#1}

\let\Algorithm\algorithm
\renewcommand\algorithm[1][]{\Algorithm[#1]\setstretch{1.5}}
%\renewcommand{\listingscaption}{Листинг}

\usepackage{pifont}
\usepackage{calc}
\usepackage{suffix}
\usepackage{csquotes}
\DeclareQuoteStyle{russian}
    {\guillemotleft}{\guillemotright}[0.025em]
    {\quotedblbase}{\textquotedblleft}
\ExecuteQuoteOptions{style=russian}
\newcommand{\enq}[1]{\enquote{#1}}  
\newcommand{\eng}[1]{\begin{english}#1\end{english}}
% Подчиненные счетчики в окружениях http://old.kpfu.ru/journals/izv_vuz/arch/sample1251.tex
\newcounter{cTheorem} 
\newcounter{cStatement}
\newcounter{cDefinition}
\newcounter{cConsequent}
\newcounter{cExample}
\newcounter{cLemma}
\newcounter{cConjecture}
\newtheorem{Theorem}{Теорема}[cTheorem]
\newtheorem{Statement}{Утверждение}[cStatement]
\newtheorem{Definition}{Определение}[cDefinition]
\newtheorem{Consequent}{Следствие}[cConsequent]
\newtheorem{Example}{Пример}[cExample]
\newtheorem{Lemma}{Лемма}[cLemma]
\newtheorem{Conjecture}{Гипотеза}[cConjecture]

\renewcommand{\theTheorem}{\arabic{Theorem}}
\renewcommand{\theDefinition}{\arabic{Definition}}
\renewcommand{\theStatement}{\arabic{Statement}}
\renewcommand{\theConsequent}{\arabic{Consequent}}
\renewcommand{\theExample}{\arabic{Example}}
\renewcommand{\theLemma}{\arabic{Lemma}}
\renewcommand{\theConjecture}{\arabic{Conjecture}}
%\makeatletter
\NewDocumentCommand{\Newline}{}{\text{\\}}
\newcommand{\sequence}[2]{\ensuremath \left(#1,\ \dots,\ #2\right)}

\definecolor{mygreen}{rgb}{0,0.6,0}
\definecolor{mygray}{rgb}{0.5,0.5,0.5}
\definecolor{mymauve}{rgb}{0.58,0,0.82}
\renewcommand{\listalgorithmname}{Список алгоритмов}
\floatname{algorithm}{Листинг}
\renewcommand{\lstlistingname}{Листинг}
\renewcommand{\thealgorithm}{\arabic{algorithm}}

\newcommand{\refAlgo}[1]{(листинг \ref{#1})}
\newcommand{\refImage}[1]{(рисунок \ref{#1})}

\renewcommand{\theenumi}{\arabic{enumi}.}% Меняем везде перечисления на цифра.цифра	
\renewcommand{\labelenumi}{\arabic{enumi}.}% Меняем везде перечисления на цифра.цифра
\renewcommand{\theenumii}{\arabic{enumii}}% Меняем везде перечисления на цифра.цифра
\renewcommand{\labelenumii}{(\arabic{enumii})}% Меняем везде перечисления на цифра.цифра
\renewcommand{\theenumiii}{\roman{enumiii}}% Меняем везде перечисления на цифра.цифра
\renewcommand{\labelenumiii}{(\roman{enumiii})}% Меняем везде перечисления на цифра.цифра
\renewcommand{\labelitemi}{---}
\renewcommand{\labelitemii}{---}

\makeatletter
\def\p@subsection{}
\def\p@subsubsection{\thesection\,\thesubsection\,}
\makeatother
\newcommand{\anonsection}[1]{\cleardoublepage
\phantomsection
\addcontentsline{toc}{section}{\protect\numberline{}#1}
\section*{#1}\vspace*{2.5ex} % По госту положены 3 пустые строки после заголовка ненумеруемого раздела
}
\newcommand{\sectionbreak}{\clearpage}
\renewcommand{\sectionfont}{\normalsize} % Сбиваем стиль оглавления в стандартный
\renewcommand{\cftsecleader}{\cftdotfill{\cftdotsep}} % Точки в оглавлении напротив разделов

\renewcommand{\cftsecfont}{\normalfont\large} % Переключение на times в содержании
\renewcommand{\cftsubsecfont}{\normalfont\large} % Переключение на times в содержании

\setlength{\cftsecindent}{0pt}% Убираем отступы в содержании для \section
\setlength{\cftsubsecindent}{0pt}% Убираем отступы в содержании для \subsection
\setlength{\cftsubsubsecindent}{0pt}% Убираем отступы в содержании для \subsubsection

\usepackage{caption} 
%\captionsetup[table]{justification=raggedleft} 
%\captionsetup[figure]{justification=centering,labelsep=endash}
\usepackage{amsmath}    % \bar    (матрицы и проч. ...)
\usepackage{amsfonts}   % \mathbb (символ для множества действительных чисел и проч. ...)
\usepackage{mathtools}  % \abs, \norm
    \DeclarePairedDelimiter\abs{\lvert}{\rvert} % операция модуля
    \DeclarePairedDelimiter\norm{\lVert}{\rVert} % операция нормы
\DeclareTextCommandDefault{\textvisiblespace}{%
  \mbox{\kern.06em\vrule \@height.3ex}%
  \vbox{\hrule \@width.3em}%
  \hbox{\vrule \@height.3ex}}    
\newsavebox{\spacebox}
\begin{lrbox}{\spacebox}
\verb*! !
\end{lrbox}
\newcommand{\aspace}{\usebox{\spacebox}}
\DeclareTotalCounter{listing}
\def\figurename{Рисунок}

\makeatletter
\renewcommand*{\p@subsubsection}{}
\makeatother

%forest
\usepackage[edges]{forest}
\usepackage{forest}

\tikzset{
  edgelabel/.style={font=\ttfamily\footnotesize, inner sep=1pt,midway, sloped, above},
  leaflabel/.style={draw=none, font=\scriptsize\ttfamily}
}



\newcommand{\bmstuheader}
{\vspace*{-30pt}
\hspace{-45pt}
\begin{minipage}{0.17\textwidth}
\hspace*{-20pt}\centering
\includegraphics[width=1.3\textwidth]{emblem.png}
\end{minipage}
\begin{minipage}{0.82\textwidth}\small \textbf{
\vspace*{-0.7ex}
\hspace*{-10pt}\centerline{Министерство науки и высшего образования Российской Федерации}
\vspace*{-0.7ex}
\centerline{Федеральное государственное автономное образовательное учреждение }
\vspace*{-0.7ex}
\centerline{высшего образования}
\vspace*{-0.7ex}
\centerline{<<Московский государственный технический университет}
\vspace*{-0.7ex}
\centerline{имени Н.Э. Баумана}
\vspace*{-0.7ex}
\centerline{(национальный исследовательский университет)>>}
\vspace*{-0.7ex}
\centerline{(МГТУ им. Н.Э. Баумана)}}
\end{minipage}

\vspace{-2pt}
\hspace{-34.5pt}\rule{\textwidth}{0.5pt}

\vspace*{-18.3pt}
\hspace{-34.5pt}\rule{\textwidth}{2.5pt}
 
\vspace{0.5ex}
\noindent \small ФАКУЛЬТЕТ\hspace{80pt} <<Информатика и системы управления>>

\vspace*{-16pt}
\hspace{35pt}\rule{0.855\textwidth}{0.4pt}

\vspace{0.5ex}
\noindent \small КАФЕДРА\hspace{50pt} <<Теоретическая информатика и компьютерные технологии>>

\vspace*{-16pt}
\hspace{25pt}\rule{0.875\textwidth}{0.4pt}
}


\begin{document}



\begin{Definition}
\label{Dict}
    \underline{Словарь} $S$ определяет язык $L$, если $L$ содержит в себе те и только те слова, которые включают в себя как подстроки все слова из $S$ (в произвольном порядке, произвольной кратности, также в слова из $L$ могут входить и другие произвольные подстроки). 
\end{Definition}

В определенных выше терминах изучим особенности задания языков через словари и, в частности, с помощью фактор-кодов --- а именно, решим следующие задачи:

\begin{enumerate}
    \item Пусть дано множество строк --- словарь $S$. Необходимо построить по нему фактор-код $F$ такой, что $L(F) = L(S)$, т.е. множества $S$ и $F$ задают один и тот же язык. 
    \item Пусть даны фактор-коды $F_1$ и $F_2$. Необходимо построить по ним эффективное предствление их объединения, т.е. фактор-код языка $L(F_1 \cup F_2 )$.
    \item Пусть даны фактор-коды $F_1$ и $F_2$. Необходимо построить минимальный фактор-код, накрывающий $L(F_1) \cup L(F_2)$. 
\end{enumerate}


Заметим, что словарь, минимальный для своего языка с точки зрения количества входящих в него элементов, будет являться фактор-кодом. Действительно, пусть дан словарь $S$, не являющийся фактор-кодом, и пусть он для его языка является минимальным. Тогда 
$$ \exists (x_1, ..., x_k, y_1, ..., y_k \in S )\; (x_1...x_k = y_1...y_k \; \& \; \exists i = 1...k \quad x_i \neq y_i). $$

Возьмем пару $x_i, \; y_i$, в которой впервые произошло несовпадение (т.е. $x_j~=~y_j \; \forall j = 1...i-1$). Так как строки $x_1...x_k$ и $y_1...y_k$ посимвольно совпадают, то $x_i$ является префиксом $y_i$ или наоборот. Пусть для определенности $y_i = x_i p$, тогда $S \; \backslash \; \{x_i\}$ задает тот же язык, что и $S$ --- это следует из определения словаря (подстрока $x_i$ автоматически содержится во всех словах языка, т.к. содержится в $y_i$). Получили противоречие к поставленному изначально утверждению --- значит исходный словарь не мог являться минимальным.

\subsection{Расшифровка}

Рассмотрим поподробнее пункты задания. 

\vspace{10pt}

{\centering
1. Построение по словарю фактор-кода.

}

\vspace{4pt}

Обратимся снова к рассуждению о минимальном словаре, приведенному ранее. Из него формально сформулируем следующие утверждения.

\begin{Statement}
\label{stat::min}
Если словарь $S$ определяет язык $L$ и при этом найдутся различные $u,\, w$ из $S$ такие, что $u$ --- подстрока $w$, то словарь $S' = S \backslash \{u\}$ также определяет язык $L$. 
\end{Statement}

Утверждение~\ref{stat::min} следует из определения словаря.

\begin{Statement}
Удалением из произвольного словаря строк, которые являются собственными подстроками других элементов словаря, можно получить фактор-код.     
\end{Statement}
\begin{proof}
    Проведем операцию удаления столько раз, сколько возможно. Предположим, что полученное множество не оказалось фактор-кодом.
    \newline
    Тогда, аналогично рассуждению в вводной части, найдутся $ x_1...x_k = y_1...y_k  $ и $x_i \neq y_i$. Из этого следует, что в множестве после удаления осталась строка, которая является подстрокой другой --- противоречие доказывает утверждение.
\end{proof}

Таким образом, для решения первой подзадачи достаточно уметь внутри некоторого множества эффективно проверять строки на включение друг в друга в качестве подстроки.

\vspace{10pt}

{\centering
2. Построение фактор-кода $L(F_1 \cup F_2)$

}

\vspace{4pt}

В данном случае множество $F_1\cup F_2$ можно рассматривать как произвольный словарь. Тогда его можно привести к виду фактор-кода так же, как в предыдущей подзадаче.

Отметим также (что неважно для дальнейшего рассуждения), что $L(F_1 \cup F_2) = L(F_1) \cap L(F_2)$ --- это верно из определения словаря (все слова языка будут включать строки, входящие и в $F_1$ и в $F_2$).


\vspace{10pt}
{\centering
3. Построение фактор-кода $L(F_1) \cup L(F_2)$

}

\vspace{4pt}

Для решения этой задачи нужно опереться на определение словаря. Пусть $S = \{w_0, w_1, ..., w_n\}$ --- словарь, тогда для каждого слова $\alpha \in L(S)$ верно, что $\alpha$ включает $w_i$ как подстроку для всех $i$.

В обратную сторону это соотношение можно построить не для всех языков. Находясь в классе задаваемых словарем языков, их объединение зачастую не будет лежать в этом же классе. В таком случае можно лишь определить наиболее узкий язык, заданный словарем, который приближен к искомому объединению.

\begin{Example}
    Пусть $L_1 = L(\{ab,\, ba,\, cc\})$ и  $L_2 = L(\{bcc,\, cab\})$.
    Тогда в словах $L_3 = L_1 \cup L_2$ гарантированно будут встречаться подстроки $a, b, c$. Более того, т.к. $ab$ содержится в словах $L_1$, а в $L_2$ --- $cab$, то в словах языка $L_3$ будет встречаться подстрока $ab$. Аналогично словарь для $L_3$ содержит $cc$. \\
    Таким образом, $dict(L_3) = \{ab, cc, a, b, c\} \sim \{ab, cc\}$. Однако, как легко видеть, $L_4 =L(\{ab, cc\}) \neq L_3$, т.к. $L_4 \supset L_3$, но $L_4$ также включает множество слов, не относящихся ни к $L_1$, ни к $L_2$.   
\end{Example}

Пусть даны два фактор-кода $F_1, \, F_2$, цель --- найти фактор-код $F_3$, в некотором смысле наиболее близкий к языку $L(F_1) \cup L(F_2)$.
Очевидно, что если некоторая строка $\omega$ содержится и в $F_1$, и в $F_2$, то $\omega \in F_3$. Более того, если $\omega$ --- подстрока некоторой $\alpha = a_1\omega a_2 \in F_1 $ и некоторой $\beta = b_1 \omega b_2 \in F_2$, то $\omega \in F_3$.

\begin{Statement}
    Пусть $F_1, \, F_2$ --- два фактор-кода. Пусть $F_3$ --- множество, каждый элемент которого --- строка, которая является общей подстрокой некоторых $\alpha$ и $\beta$, причем $\alpha \in F_1 ,\; \beta \in F_2$. 
    \\
    Тогда $L(F_1) \cup L(F_2) \, \subset L(F_3)$.
\end{Statement}

\begin{proof}
    Проверим $L(F_1) \subset L(F_3)$ и $L(F_2) \subset L(F_3)$.
    Пусть произвольное слово $\omega \in L(F_1), \; F_1 = \{ u_1, u_2, ..., u_n\}, \; 
    F_3 = \{p_1, p_2, ..., p_k\}$.
    \\
    При этом верно, что $u_i$ --- подстрока $\omega$ для всех $i$. Тогда верно, что $p_i$ --- подстрока $\omega$ для всех $i$ --- это автоматически выполняется, т.к. $p_i$ --- подстрока некоторого~$u_j$. 
    Так как любое слово языка $L(F_1)$ также задается словарем $F_3$, то 
    $L(F_1) \subset L(F_3)$. Аналогичное соотношение верно для $L(F_2)$.
\end{proof}

\begin{Statement}
    Определенный выше фактор-код $F_3$ является "наиболее точным" \, покрытием
    $L(F_1) \cup L(F_2) $
    в том смысле, что 
    $$ \nexists F_4\ | \ L(F_1) \cup L(F_2) \, \subset L(F_4) \, \subset L(F_3)$$.
\end{Statement}

\begin{proof}
    Предположим обратное. Тогда, принимая во внимание, что $L(F_4) \neq L(F_3)$, верно, что $\exists \omega | \; \omega \in L(F_3) \; \& \; \omega \notin L(F_4)$. Из определения словаря получим, что $\exists u \in F_4 | \; \omega \not \succ u$. С другой стороны, из соотношения
    $ L(F_1) \cup L(F_2) \, \subset L(F_4) $ получим, что строка $u$ содержится как подстрока во всех словах языков $L(F_1)$ и $L(F_2)$.
    \\
    Если $u$ не является подстрокой  некоторых строк из фактор-кодов $F_1$ и $F_2$, то невозможно гарантировать, что $u$ содержится во всех словах $L(F_1)$ и $L(F_2)$. Предположим, что символ $\textbf{x}$ не встречается в $u$. Тогда слово $\alpha \in L(F_1), \; \alpha = s_1\textbf{x}s_2...\textbf{x}s_k$ ($s_i$ --- элементы $F_1$) не может содержать $u$. Аналогичное рассуждение верно для $L(F_2)$. Т.к. предполагается, что алфавит бесконечно большой, символ $\textbf{x}$ всегда найдется.
    \\
    Отсюда получаем противоречие, т.к. строка $u \in F_4$ должна быть подстрокой некоторых элементов фактор-кодов $F_1$ и $F_2$, но тогда эта строка содержится в фактор-коде $F_3$.
\end{proof}


Таким образом, для решения третьей подзадачи необходимо уметь эффективно находить общие подстроки.



\subsection{Агротехнологии и другие приемы}

Конечно, во всем перечисленном выше полезно уметь сажать деревья. 

Построив эффективно обобщенное суффиксное дерево для двух множеств, мы можем за время, зависящее от длины строки, проверять ее включение как подстроки, добавлять строку в дерево и находить максимальные общие подстроки.






\begin{figure}
    \centering
    \begin{forest}
for tree={
  l sep=24mm,                  % level separation
  s sep=15mm,                   % sibling separation
  draw,
  circle,
  minimum size=2.2pt,          % tiny dot nodes
  inner sep=0pt,
  edge={-Latex},               % arrows on edges
}
[root, leaflabel % root
  % b*
  [, edge label={node[edgelabel]{ \$ a n a n a b}} ]

  % $ (the standalone suffix)
  [10 5, leaflabel, edge label={node[edgelabel, sloped, above]{\$}} ]

  % a*
  [, edge label={node[edgelabel, sloped, above]{a}}
    % a$
    [, edge label={node[edgelabel,above]{\$}} ]
    % a + na*
    [, edge label={node[edgelabel,above]{n a}}
      % a + na + $
      [, edge label={node[edgelabel,above]{\$}} ]
      % a + na + na$
      [, edge label={node[edgelabel,sloped, above]{n a \$}} ]
    ]
  ]

  % n*
  [, edge label={node[edgelabel,above]{n a}}
    % n + a + $
    [, edge label={node[edgelabel,above]{\$}} ]
    % n + a + na$
    [, edge label={node[edgelabel,above]{n a \$}} ]
  ]
]
\end{forest}

    \caption{Caption}

    \label{fig:placeholder}

\end{figure}

\begin{algorithm}


\end{algorithm}



\end{document}
